\section{Cycle Integral Definitions}

The cycle integral extends the finite integral to unbounded repeating
sequences. Two equivalent definitions are proven.

\begin{itemize}
  \item Recursive: recurrence on the cycle position --- Subsection~3.1
  \item Modulo: closed-form using \texttt{div} and \texttt{mod} ---
    Subsection~3.2
  \item The two definitions are extensionally equivalent --- Subsection~3.3
\end{itemize}

\begin{equation*}
\forall\, i \in \mathbb{N}_0,\ init \in \mathbb{N}_0,\ L \in \mathbb{L},\ n = |L|
\end{equation*}

\begin{equation*}
\operatorname{CycleIntegral}(L, init) := [w_0, w_1, w_2, \ldots]
\end{equation*}

\begin{equation*}
\begin{aligned}
0 \leq i < n \implies w_i = \sum_{j=0}^i L_j + init \quad &\text{[Sum Property]} \\
i > 0 \implies  \ w_i - w_{i-1} = L_{(i \operatorname{mod} n)}
\quad &\text{[Step Property]}
\end{aligned}
\end{equation*}

\subsection{Recursive Cycle Integral}

\begin{equation*}
\begin{aligned}
\operatorname{RecCycle}(L) &= [v_0, v_1, \dots] \mid v_i =
\begin{cases}
  L[i] & i < n \\
  v_{i - n} & i \geq n
\end{cases} \\
\operatorname{RecCycleIntegral}(L, init) &= [w_0, w_1, \dots] \mid w_i =
\begin{cases}
  init + w_0 & i = 0 \\
  v_i + w_{i - 1} & i > 0
\end{cases}
\end{aligned}
\end{equation*}

\textbf{Recursive Cycle Equivalence}: Recursive Cycle is equivalent to Mod
Cycle, as proven in the article
\href{https://github.com/thiagomata/prime-numbers/blob/master/articles/chapter4/cycle.md}{\emph{Formal
Verification of Cyclic Lists}}~\cite{mata2026cycle}.

\begin{equation*}
\begin{aligned}
\operatorname{RecCycle}(L)_i &=  \operatorname{ModCycle}(L)_i \quad &\text{[Cycle Equivalence]} \\
  &= L_{(i \operatorname{mod} n)} \quad &\text{[Mod Cycle Definition]}
\end{aligned}
\end{equation*}

\textbf{Sum Property}:

\begin{equation*}
\begin{aligned}
v_0  &= L_0  \quad &\text{[Base Case]} \\
&= L_{(0 \operatorname{mod} n)} \quad &\text{[By Modulo Property]} \\
&= \sum_{j=0}^0 L_j \quad &\text{[Summation Re-indexing]}
\end{aligned}
\end{equation*}

\begin{equation*}
\begin{aligned}
& i < n \implies \\
w_0 &= init + v_0 \quad &\text{[Base Case]} \\
&= init + \sum_{j=0}^i L_j  \quad &\text{[Summation Re-indexing]} \\
& i > 0 \implies \\
w_i &= init + v_0 + \sum_{j=1}^i v_j \quad &\text{[By Definition]}\\
&= init + v_0 + \sum_{j=1}^i L_j \quad &\text{[Mod of Small Values Property]} \\
&= init + \sum_{j=0}^i v_j \quad &\text{[Summation Re-indexing]} \\
\therefore w_i &= init + \sum_{j=0}^i L_j \quad &\text{[Q.E.D.]}
\end{aligned}
\end{equation*}

{\raggedright
This property is verified in the
\href{https://github.com/thiagomata/prime-numbers/blob/master/src/main/scala/v1/chapter4/cycle/integral/recursive/properties/CycleIntegralProperties.scala}{\texttt{CycleIntegralProperties::\allowbreak assertCycleIntegralEqualsSumSmallPositions}}.
A key Scala verification excerpt is in Appendix~A.1; the complete proof is
linked in the source reference.\par}

\textbf{Step Property}:

\begin{equation*}
\begin{aligned}
w_i - w_{i-1} &= v_i + w_{i-1} - w_{i-1} \quad &\text{[By Definition]} \\
&= v_i \quad &\text{[Simplification]} \\
&= L_{(i \operatorname{mod} n)} \quad &\text{[Substitution]}
\end{aligned}
\end{equation*}

{\raggedright
This property is verified in the
\href{https://github.com/thiagomata/prime-numbers/blob/master/src/main/scala/v1/chapter4/cycle/integral/recursive/properties/CycleIntegralProperties.scala}{\texttt{CycleIntegralProperties::\allowbreak assertDiffEqualsCycleValue}}.
A key Scala verification excerpt is in Appendix~A.2; the complete proof is
linked in the source reference.\par}

\subsection{Modulo Cycle Integral}

\begin{equation*}
\begin{aligned}
&\operatorname{ModCycle}(L)_i &:= L_{(i \operatorname{mod} n)} = [w_0, w_1, \dots ] \\
&I_k &:= \sum_{j=0}^{k} L_j \quad (0 \leq k < n) \quad &\text{[Integral of L]} \\
&S &:= I_{n-1} \quad &\text{[One full cycle sum]} \\
&\operatorname{ModCycleIntegral}(L, init)_i &:= (i \operatorname{div} n)\cdot S + I_{(i \operatorname{mod} n)} + init
\end{aligned}
\end{equation*}

\textbf{Sum Property}:

\begin{equation*}
i < n \implies w_i = \sum_{j=0}^i L_j + init \quad \text{[Claim to Prove]}
\end{equation*}

\begin{equation*}
\begin{aligned}
& i < n \implies  \\
& i \operatorname{div} n
      \quad = 0 \quad
      \text{[By Div of Small Values Property]} \\
w_i &= (i \operatorname{div} n)\cdot S + I_{(i \operatorname{mod} n)} + init \quad
      &\text{[Definition]} \\
&= 0 \cdot S + I_i + init \quad &\text{[Substitution]} \\
&= \sum_{j=0}^i L_j + init \quad &\text{[By definition of } I_i]
\end{aligned}
\end{equation*}

\begin{equation*}
\therefore \
\forall\, i < n,\quad w_i = \sum_{j=0}^i L_j + init \quad \text{[Q.E.D.]}
\end{equation*}

{\raggedright
This property is verified in the
\href{https://github.com/thiagomata/prime-numbers/blob/master/src/main/scala/v1/chapter4/cycle/integral/mod/ModCycleIntegralProperties.scala}{\texttt{ModCycleIntegralProperties::\allowbreak assertFirstValuesMatchIntegral}}.
A key Scala verification excerpt is in Appendix~A.3; the complete proof is
linked in the source reference.\par}

\textbf{Step Property}:

\begin{equation*}
w_i - w_{i-1} = L_{\, i \operatorname{mod} n}, \quad i>0,\, n>0
\quad \text{[Claim to Prove]}
\end{equation*}

\textbf{Case $i \operatorname{mod} n > 0$:}

\begin{equation*}
\begin{aligned}
&i \operatorname{mod} n &= ((i-1) \operatorname{mod} n) + 1 &&\text{[By Modulo Properties]} \\
&i \operatorname{div} n &= (i-1) \operatorname{div} n &&\text{[By Division Properties]} \\
&w_i &= (i \operatorname{div} n)\,S + I_{\,i\operatorname{mod} n} + init &&\text{[Definition]} \\
     &&= (i-1 \operatorname{div} n)\,S + I_{\,((i-1)\operatorname{mod} n)+1} + init &&\text{[Div/Mod property]} \\
&w_{i-1} &= (i-1 \operatorname{div} n)\,S + I_{\, (i-1)\operatorname{mod} n} + init &&\text{[Definition]} \\
&w_i-w_{i-1} &= I_{\,((i-1)\operatorname{mod} n)+1} - I_{\, (i-1)\operatorname{mod} n} &&\text{[Cancellation]} \\
    &&= \Big(\sum_{j=0}^{(i-1)\operatorname{mod} n} L_j + L_{((i-1)\operatorname{mod} n)+1}\Big)
       - \sum_{j=0}^{(i-1)\operatorname{mod} n} L_j &&\text{[Expand sum]} \\
    &&= L_{((i-1)\operatorname{mod} n)+1} &&\text{[Cancellation]} \\
    &&= L_{\,i \operatorname{mod} n} &&\text{[Modulo property]}.
\end{aligned}
\end{equation*}

\textbf{Case $i \operatorname{mod} n = 0$:}

\begin{equation*}
\begin{aligned}
w_i &= (i \operatorname{div} n)\,S + L_0 + init
&&\text{[Definition]} \\
w_{i-1} &= (i \operatorname{div} n -1)\,S + I_{\,n-1} + init
&&\text{[Div/Mod property]} \\
w_i-w_{i-1}
    &= (i \operatorname{div} n)\,S + L_0 + init
       - \big((i \operatorname{div} n -1)\,S + I_{\,n-1} + init\big)
&&\text{[Substitution]} \\
    &= S + L_0 - I_{\,n-1}
&&\text{[Simplification]} \\
    &= L_0
&&\text{[Since } S = I_{\,n-1}] \\
    &= L_{\,i \operatorname{mod} n}
&&\text{[Modulo property]}.
\end{aligned}
\end{equation*}

\begin{equation*}
\therefore \
w_i - w_{i-1} = L_{\, i \operatorname{mod} n}, \quad \forall\, i > 0 \quad \text{[Q.E.D.]}
\end{equation*}

{\raggedright
This property is verified in the
\href{https://github.com/thiagomata/prime-numbers/blob/master/src/main/scala/v1/chapter4/cycle/integral/mod/ModCycleIntegralProperties.scala}{\texttt{ModCycleIntegralProperties::\allowbreak assertSimplifiedDiffValuesMatchCycle}}.
A key Scala verification excerpt is in Appendix~A.4; the complete proof is
linked in the source reference.\par}

\subsection{Equivalence of Definitions}

The nontrivial equivalence proved here is between the recursive
\texttt{CycleIntegral} definition and the closed-form
\texttt{ModCycleIntegral} definition.

\begin{equation*}
\begin{aligned}
\operatorname{CycleIntegral}(L, init)
&= \operatorname{ModCycleIntegral}(L, init) \\
&= [w_0, w_1, w_2, \ldots] \mid w_i = \sum_{j=0}^i L_{(j \operatorname{mod} n)} + init \ \blacksquare
\end{aligned}
\end{equation*}

{\raggedright
This property is verified in the
\href{https://github.com/thiagomata/prime-numbers/blob/master/src/main/scala/v1/chapter4/cycle/integral/mod/ModCycleIntegralProperties.scala}{\texttt{ModCycleIntegralProperties::\allowbreak assertCycleIntegralMatchModCycleDef}}.
A key Scala verification excerpt is in Appendix~A.5; the complete proof is
linked in the source reference.\par}
